As discussed in \Cref{sec:mathematical_underpinnings,sec:trapdoor_functions}, cryptography is ultimately based on mathematics. Therefore, the theoretical design and logical robustness of cryptographic algorithms and protocols are often scrutinized to reveal any weaknesses or attacks.

However, the \textbf{application of cryptography} in real-world scenarios yields security concerns other than only those related to its mathematical underpinnings. Even more, cryptography in general is a necessary but insufficient condition for guaranteeing the security of sensitive data and services \cite{menezes_challenges_2021}. In fact, \textbf{(applied) cryptography may break for a variety of reasons} not strictly related to mathematics: implementation errors and vulnerabilities, (unintentionally introduced) backdoors, misconfigurations and misuses, wrong security assumptions or standard definitions, side-channel attacks, and social engineering attacks, to mention a few.

Below, we provide an overview of security in cryptography starting from its most generic aspects: the definition of a threat model (\Cref{sec:cryptographic_security_overview:threat_model}), the consequent choice of attack model (\Cref{sec:cryptographic_security_overview:attack_model}), and a taxonomy of attacks (\Cref{sec:cryptographic_security_overview:attacks}).